\chapter{Model-based search: Bayesian optimization and its relatives}
\label{sec:bayesian}

Random search treats trial fifty exactly like trial one: it has learned
nothing. When each trial costs GPU-hours, that indifference is expensive,
and it pays to think hard between trials. Bayesian optimization (BO) spends
that thought in two steps: fit a probabilistic model of the objective from
the trials seen so far (the \emph{surrogate}), then choose the next
configuration by maximizing an \emph{acquisition function} that scores
candidate points by how useful evaluating them would be. The surrogate is
cheap to evaluate, so the inner maximization can afford thousands of
candidate evaluations to decide where to spend one real trial. Modern
references are the book by \citet{garnett2023bayesoptbook} and the survey by
\citet{shahriari2016taking}; here we develop the machinery far enough that
the design choices in our product are not acts of faith, starting from the
question the surrogate must answer: after eight trials of the running
example, what do we believe about the loss at a learning rate we have never
tried, and how sure are we?

\section{Gaussian process surrogates}
\label{sec:gp}

The default surrogate is a Gaussian process (GP), and its entire content is
one modeling assumption stated positively: we model the objective $f$ as a
random function such that any finite collection of values
$(f(\config_1), \dots, f(\config_n))$ is jointly Gaussian, with prior mean
zero and covariance given by a kernel,
$\operatorname{Cov}[f(\config), f(\config')] = k(\config, \config')$. The
kernel is where all the prior knowledge lives: it says how similar the
objective's values at two configurations should be as a function of how far
apart the configurations are. The common default, and ours, is the
Mat\'ern-5/2 kernel with a lengthscale per dimension,
\begin{equation}
k(\config, \config')
= \sigma_f^2 \left( 1 + \sqrt{5}\,\rho + \tfrac{5}{3}\rho^2 \right)
  e^{-\sqrt{5}\,\rho},
\qquad
\rho^2 = \sum_{j=1}^{d} \frac{(\lambda_j - \lambda'_j)^2}{\ell_j^2},
\label{eq:matern}
\end{equation}
which produces functions that are twice differentiable but not infinitely
smooth, a better match for loss surfaces than the very smooth squared
exponential \citep{rasmussen2006gaussian}. Each lengthscale $\ell_j$ has a
direct reading: it is the distance along axis $j$ over which the objective's
values decorrelate. A long lengthscale on the entropy-coefficient axis says
that moving the entropy coefficient barely changes the loss.

It is worth pausing on why the model is nonparametric, because the
obvious alternative is one every applied statistician has fit. Response
surface methodology would posit a global parametric form, a quadratic in
the knobs, and estimate its coefficients by least squares. That imposes
one curvature on the entire space: the model must believe the loss bends
the same way near $10^{-5}$ as near $10^{-2}$, and eight trials estimating
a fifteen-dimensional quadratic are hopeless anyway (the quadratic alone
has over a hundred coefficients). The kernel assumption is weaker and
better matched to what we actually know: not the shape of the surface,
only that configurations close together perform similarly, with
``close'' learned per axis from the data. The GP is what Bayesian
regression becomes when the parametric form is dropped and only that
smoothness belief is kept; the price is that predictions lean on nearby
data rather than on a global formula, which is exactly the right
epistemic posture for a surface we know nothing else about.

We assume observations are the true objective plus Gaussian noise,
$y_i = f(\config_i) + \varepsilon_i$ with
$\varepsilon_i \sim \N(0, \sigma^2)$, which is our stand-in for seed
variation; it is a modeling choice, adopted because it keeps every formula
below in closed form, and Section~\ref{sec:gp-practical} returns to what to
do when the noise is visibly not Gaussian. Under these assumptions the
prediction problem solves itself by pure linear algebra. The observed vector
$\mathbf{y} = (y_1, \dots, y_n)$ and the unknown value $f(\config)$ at a new
point are jointly Gaussian,
\begin{equation}
\begin{pmatrix} \mathbf{y} \\ f(\config) \end{pmatrix}
\sim \N\!\left(
\mathbf{0},
\begin{pmatrix}
\mathbf{K} + \sigma^2 \mathbf{I} & \mathbf{k}(\config) \\
\mathbf{k}(\config)^\top & k(\config, \config)
\end{pmatrix}
\right),
\qquad
\mathbf{K}_{ij} = k(\config_i, \config_j),
\quad
\mathbf{k}(\config)_i = k(\config_i, \config),
\label{eq:joint}
\end{equation}
and the standard formula for conditioning one block of a Gaussian on
another (condition the second block on the first) gives the posterior at
$\config$ as another Gaussian,
\begin{equation}
\mu_n(\config) = \mathbf{k}(\config)^\top
\left( \mathbf{K} + \sigma^2 \mathbf{I} \right)^{-1} \mathbf{y},
\qquad
s_n^2(\config) = k(\config, \config) - \mathbf{k}(\config)^\top
\left( \mathbf{K} + \sigma^2 \mathbf{I} \right)^{-1} \mathbf{k}(\config)
\label{eq:gp}
\end{equation}
\citep{rasmussen2006gaussian}. Readers who know kriging from spatial
statistics or Bayesian linear regression from econometrics are on home
ground here: \eqref{eq:gp} is exactly that conditioning, with the kernel
playing the role of the spatial covariance function.

The two formulas
carry the whole story, and they are worth reading term by term. The posterior mean is a weighted
combination of the observed values, with weights
$\mathbf{k}(\config)^\top (\mathbf{K} + \sigma^2 \mathbf{I})^{-1}$ set by
similarity under the kernel: observations near $\config$ (as measured in
lengthscale units) dominate the prediction.

The posterior variance starts
from the prior variance $k(\config,\config)$ and \emph{subtracts} a term
that grows as data approaches $\config$; at fixed kernel
hyperparameters nothing about the observed \emph{values} appears in it,
only the observed \emph{locations}, so the model knows where it is
ignorant before it knows what the answer is. The qualifier is not
decoration: once the lengthscales and the noise variance are estimated
from the data by \eqref{eq:mll}, the observed values do influence the
variance, through the fitted hyperparameters rather than through the
conditioning formula.

The noise term $\sigma^2 \mathbf{I}$ deserves its own sentence: with
$\sigma^2 = 0$ the posterior mean passes exactly through every
observation, treating each score as gospel, while $\sigma^2 > 0$ turns
interpolation into smoothing, letting the mean pass \emph{near} noisy
scores instead of through them. Readers who know ridge regression will
recognize the mechanism, since $(\mathbf{K} + \sigma^2\mathbf{I})^{-1}$
applies precisely ridge's shrinkage, with the noise variance as the
regularization constant; estimating $\sigma^2$ from the data by
\eqref{eq:mll} below is what lets the model decide for itself how much
of the spread in RL returns is signal and how much is seed luck.

Figure~\ref{fig:gp} shows the shape this produces on the running example's
learning-rate axis: the band of uncertainty is pinched near the trials we
have paid for and swells between and beyond them, which is exactly the
information an acquisition function needs.

\begin{figure}[t]
\centering
\begin{tikzpicture}[scale=1.0, font=\small]
% axes
\draw[-{Latex}] (0,0) -- (9,0);
\node[anchor=north] at (7.4,-0.12) {learning rate (log scale)};
\draw[-{Latex}] (0,0) -- (0,4.4);
\node[rotate=90, anchor=south] at (-0.4,2.2) {validation loss};
% uncertainty band (upper then lower reversed)
\fill[blue!12]
  (0.3,3.9) .. controls (1.0,2.9) and (1.4,2.55) .. (1.8,2.45)
  .. controls (2.6,2.25) and (3.2,1.75) .. (3.9,1.55)
  .. controls (4.6,1.35) and (5.2,1.30) .. (5.8,1.45)
  .. controls (6.6,1.70) and (7.4,2.9) .. (8.4,3.9)
  -- (8.4,2.1)
  .. controls (7.4,1.6) and (6.6,1.15) .. (5.8,1.05)
  .. controls (5.2,0.95) and (4.6,1.00) .. (3.9,1.15)
  .. controls (3.2,1.35) and (2.6,1.85) .. (1.8,2.15)
  .. controls (1.4,2.25) and (1.0,2.3) .. (0.3,2.5)
  -- cycle;
% mean
\draw[thick, blue!60!black]
  (0.3,3.2) .. controls (1.0,2.6) and (1.4,2.4) .. (1.8,2.3)
  .. controls (2.6,2.05) and (3.2,1.55) .. (3.9,1.35)
  .. controls (4.6,1.18) and (5.2,1.12) .. (5.8,1.25)
  .. controls (6.6,1.45) and (7.4,2.2) .. (8.4,3.0);
% data points
\fill (1.8,2.3) circle (2.2pt);
\fill (3.9,1.35) circle (2.2pt);
\fill (5.8,1.25) circle (2.2pt);
\fill (7.6,2.45) circle (2.2pt);
% annotations
\draw[-{Latex}] (3.15,0.72) -- (3.8,1.12);
\node[anchor=north, align=center] at (1.85,0.78) {band pinched where\\we have paid for data};
\draw[-{Latex}] (2.9,3.4) -- (2.9,2.35);
\node[anchor=south, align=center] at (2.9,3.4) {band swells between\\observations};
\end{tikzpicture}
\caption{A Gaussian-process posterior on the running example's
learning-rate axis after four trials (schematic illustration of the
mechanism, not measured data). The mean~\eqref{eq:gp} passes near the
observations; the shaded one-standard-deviation band shows
$\mu_n \pm s_n$, which is pinched near observed points and grows away
from them. The figure is drawn in the low-noise regime for legibility:
with $\sigma^2 > 0$ the mean does not interpolate exactly and the
variance at an observed location stays strictly positive, and only in
the noiseless limit $\sigma^2 \to 0$ does the band close completely.
Given the kernel hyperparameters, the posterior variance depends only on
\emph{where} we have evaluated, not on the values we saw, so the model
maps its own ignorance; once those hyperparameters are themselves fitted
by marginal likelihood (below), the observed values re-enter the
variance indirectly through them.}
\label{fig:gp}
\end{figure}

The kernel hyperparameters (the lengthscales $\ell_j$, signal variance
$\sigma_f^2$, and noise $\sigma^2$) are learned by maximizing the marginal
likelihood of the data,
\begin{equation}
\log p(\mathbf{y} \mid \boldsymbol{\ell}, \sigma_f, \sigma)
= -\tfrac{1}{2} \mathbf{y}^\top (\mathbf{K} + \sigma^2 \mathbf{I})^{-1}
   \mathbf{y}
  \;-\; \tfrac{1}{2} \log \det (\mathbf{K} + \sigma^2 \mathbf{I})
  \;-\; \tfrac{n}{2} \log 2\pi ,
\label{eq:mll}
\end{equation}
whose two data-dependent terms pull in opposite directions: the first
rewards kernels under which the observed values are unsurprising (fit), the
second penalizes kernels that could explain anything (complexity), so
maximizing the sum performs a built-in Occam trade-off
\citep{rasmussen2006gaussian}.

This is how a GP quietly performs the
sensitivity analysis that grid search cannot: if the data can be explained
with a very long lengthscale on some axis, the likelihood prefers it, and
that axis is effectively switched off. The fitted numbers are worth
reading, not just using: if the running example's model learns a
lengthscale of half a decade on the log learning-rate axis, then trials
a decade apart are nearly independent evidence and the axis needs dense
coverage; a fitted lengthscale of five decades says the axis is
effectively flat and the search can spend its trials elsewhere.

The $n \times n$ solve in
\eqref{eq:gp} and \eqref{eq:mll} costs $O(n^3)$ time, irrelevant at our
trial counts (tens to low hundreds) but a real constraint for methods that
model learning curves point by point, a fact that returns in
Chapter~\ref{sec:multifidelity}.

\subsection{What actually goes wrong in practice}
\label{sec:gp-practical}

Two practical amendments matter more than any exotic kernel choice, and
both are the same move we made for random search: respect the data's
geometry. First, inputs should be warped: search axes like learning rate
live on a log scale, and beyond fixed transforms it helps to learn a
monotone warping of each axis (HEBO uses a Kumaraswamy distribution
function per axis) so that the objective looks stationary to the kernel;
a kernel with one lengthscale per axis assumes the objective wiggles at
the same rate everywhere along that axis, and warping is how we make that
assumption true rather than hoping it holds.

Second, outputs should be
transformed toward Gaussianity by a power transform, because raw losses are
often heavy-tailed and heteroscedastic: one diverged trial with a loss of
$10^{6}$ will otherwise dominate the fit of \eqref{eq:mll} and flatten the
model's view of the interesting region.

HEBO's authors open their paper
with hypothesis tests on 108 real tuning tasks showing that
heteroscedasticity and non-stationarity are the rule, not the exception,
and HEBO won the NeurIPS 2020 black-box optimization challenge with exactly
these amendments plus one more hedge: instead of trusting any single
rule for choosing the next trial, HEBO scores candidates by three
standard such rules at once (the acquisition functions of the next
section) and picks candidates that no rival beats on all three
\citep{cowenrivers2022hebo, turner2021bayesian}. The lesson for
implementers is encouraging: the wins come from respecting the data's
geometry, not from acquisition novelty.

\section{Acquisition functions}
\label{sec:acquisition}

The surrogate summarizes belief; the acquisition function turns belief into
a decision, and the decision has a tension in it. Sampling where the mean
$\mu_n$ is best exploits what we know; sampling where the variance $s_n^2$
is largest explores what we do not. Every acquisition function is a
different exchange rate between the two.

\subsection{Expected improvement, derived and then broken}

Expected improvement (EI), the traditional default, scores a candidate by
the expected amount by which it beats the best value seen so far. Write
$f^{\star}$ for the incumbent value (in minimization). If we evaluate at
$\config$ and observe $f(\config)$, the improvement is
$\max(f^{\star} - f(\config), 0)$: the amount by which we beat the
incumbent, or zero if we fail to. An economist will recognize the shape
immediately: this is an option payoff with strike $f^{\star}$, and EI is
its expected value under the posterior, so Bayesian optimization spends
its trials the way one prices options, paying for upside while the
downside is capped at the cost of the trial.

Under the posterior \eqref{eq:gp},
$f(\config) \sim \N(\mu_n, s_n^2)$, and the expectation of the improvement
has a closed form. Substituting $z = (f - \mu_n)/s_n$ and writing
$u = (f^{\star} - \mu_n)/s_n$ for the incumbent's position in posterior
standard units,
\begin{equation}
\alpha_{\mathrm{EI}}(\config)
= \int_{-\infty}^{f^{\star}} (f^{\star} - f)\,
  \frac{1}{s_n}\varphi\!\Big(\frac{f-\mu_n}{s_n}\Big)\, df
= s_n \int_{-\infty}^{u} (u - z)\, \varphi(z)\, dz
= s_n \left[ u\,\Phi(u) + \varphi(u) \right],
\label{eq:ei-derive}
\end{equation}
using $\int_{-\infty}^{u} \varphi(z)\,dz = \Phi(u)$ and
$\int_{-\infty}^{u} z\,\varphi(z)\,dz = -\varphi(u)$, the latter because
$\varphi'(z) = -z\varphi(z)$. Rearranged into its usual form,
\begin{equation}
\alpha_{\mathrm{EI}}(\config)
= \left( f^{\star} - \mu_n(\config) \right) \Phi(u)
  + s_n(\config)\, \varphi(u),
\qquad
u = \frac{f^{\star} - \mu_n(\config)}{s_n(\config)},
\label{eq:ei}
\end{equation}
with $\varphi$ and $\Phi$ the standard normal density and distribution
function \citep{jones1998efficient}. The two terms are the two halves of
the tension made explicit: the first rewards a promising mean (how much
better than the incumbent we expect to be, times the probability of being
better at all), the second rewards uncertainty, and the balance between
them is automatic, with no exploration knob to tune.

A small worked case shows the balance operating, and it is worth doing
once by hand. Suppose the incumbent loss in the running example is
$f^{\star} = 0.60$. Candidate A sits in unexplored territory: the
posterior says $\mu_n = 0.62$, $s_n = 0.05$, so $u = -0.4$, and
\eqref{eq:ei} gives
$\alpha_{\mathrm{EI}} = (-0.02)(0.345) + (0.05)(0.368) \approx 0.012$.
Candidate B sits next to the incumbent: $\mu_n = 0.59$, $s_n = 0.005$, so
$u = 2$, and
$\alpha_{\mathrm{EI}} = (0.01)(0.977) + (0.005)(0.054) \approx 0.010$.

Candidate A has the \emph{worse} expected loss, yet EI ranks it slightly
ahead, because its wide posterior leaves real probability of a large
improvement, while candidate B can improve on the incumbent by at most a
hair. That preference for informative gambles over safe pennies is the
entire personality of EI, and Figure~\ref{fig:ei} draws where it comes
from: the acquisition is the mass of the posterior lying below the
incumbent line, weighted by how far below it lies.

\begin{figure}[t]
\centering
\begin{tikzpicture}[scale=1.0, font=\small]
\draw[-{Latex}] (0,0) -- (9,0) node[below left=1mm and 0mm] {learning rate};
\draw[-{Latex}] (0,0) -- (0,4.4) node[rotate=90, above left=2mm and -14mm] {validation loss};
% incumbent line
\draw[dashed, thick] (0,1.6) -- (9,1.6) node[pos=0.03, above, font=\small] {incumbent $f^{\star}$};
% posterior mean curve (light)
\draw[blue!50]
  (0.4,3.0) .. controls (1.6,2.2) and (2.8,1.9) .. (4.0,2.05)
  .. controls (5.2,2.2) and (6.4,2.5) .. (8.4,3.1);
% candidate location
\draw[dotted] (4.0,0) node[below] {candidate $\config$} -- (4.0,4.1);
% rotated gaussian at x=4 centered at mu=2.05, drawn as profile to the right
\draw[thick]
  (4.0,3.55) .. controls (4.15,3.35) and (4.55,2.75) .. (4.62,2.05)
  .. controls (4.55,1.35) and (4.15,0.75) .. (4.0,0.55);
% shade improvement region: part of bell below y=1.6
\fill[blue!25]
  (4.0,1.6) -- (4.45,1.6)
  .. controls (4.3,1.15) and (4.15,0.85) .. (4.0,0.55) -- cycle;
\node[anchor=west, align=left] at (5.1,0.9)
  {shaded mass: outcomes that beat $f^{\star}$;\\
   EI weights each by its depth below the line};
\draw[-{Latex}] (5.05,0.95) -- (4.32,1.2);
\fill (4.0,2.05) circle (1.6pt);
\node[anchor=east] at (3.95,2.1) {$\mu_n(\config)$};
\end{tikzpicture}
\caption{What expected improvement measures (schematic). At the candidate,
the posterior over the unknown loss is the sideways bell curve. Only the
shaded tail below the incumbent line counts as improvement, and
\eqref{eq:ei} is exactly this tail's mass weighted by its depth. When the
whole bell sits far above the line, the shaded region, and with it EI and
its gradient, vanishes; that observation is the doorway to the numerical
failure discussed next.}
\label{fig:ei}
\end{figure}

Now the flaw, which is not conceptual but floating-point, and which
matters because we will inherit it the day we ship a naive
implementation. Once a decent incumbent exists, most of the search space
has $\mu_n$ far above $f^{\star}$, so $u$ is a large negative number. For
large negative $u$ the standard tail expansion of $\Phi$ gives
$\alpha_{\mathrm{EI}} \approx s_n \varphi(u)/u^2$, and $\varphi(u)$ decays
like $e^{-u^2/2}$: at $u = -6$ the acquisition is already of order
$10^{-10} s_n$, and near $u \approx -38$ the density $\varphi(u)$ leaves
the range of double-precision numbers entirely, so \eqref{eq:ei} evaluates
to exactly zero, its gradient with it. A value of $u = -6$ is not exotic;
it just means the posterior mean is six posterior standard deviations
worse than the incumbent, which describes most of the domain once the
search has found anything good. In that regime multi-restart gradient
ascent on the acquisition degenerates into random search among its own
starting points.

\citet{ament2023logei} analyze this failure, show that
the affected region grows with observations, dimensions, and constraints,
and repair it by optimizing numerically careful $\log$ transforms of the
whole EI family (analytic, batch, constrained, and hypervolume variants):
$\log \alpha_{\mathrm{EI}}$ is computed directly from stable asymptotic
expansions rather than by taking the logarithm of an underflowed number,
and smooth approximations replace the hard maximum inside the Monte
Carlo versions, the sampling-based batch variants that arrive later in
this section. The logarithm does not change where the maxima are; it changes
whether an optimizer can find them, because gradients of order
$10^{-300}$ become gradients of order one.

The repaired acquisitions
match or beat more elaborate alternatives across standard benchmarks,
which suggests much of EI's mixed reputation was numerics all along. Our
implementation should treat the LogEI variants as the default and never
optimize raw EI.

\subsection{The rest of the family, and when we need it}

A different family spends its budget on information rather than
improvement: predictive entropy search values a candidate by how much its
observation would shrink uncertainty about the optimum's location
\citep{hernandez2014pes}, max-value entropy search replaces the location
(a $d$-dimensional object) with the optimum's value (a scalar), making the
computation far cheaper \citep{wang2017mes}, and joint entropy search
conditions on fantasized optimal input-output pairs to capture both at
once \citep{hvarfner2022jes}.

These acquisitions are elegant and
occasionally worth their cost on very expensive problems, but given the
LogEI result above we treat them as optional extras rather than defaults;
\citet{frazier2018tutorial} is the standard tour of the whole acquisition
zoo, including the knowledge gradient that Chapter~\ref{sec:moo} meets
again in multi-objective form.

Two families beyond EI matter for us directly. Upper confidence bound
(UCB) style rules score a candidate by an optimistic quantile of the
posterior,
\begin{equation}
\alpha_{\mathrm{UCB}}(\config)
= \mu_n(\config) - \beta^{1/2}\, s_n(\config)
\label{eq:ucb}
\end{equation}
for minimization: pretend every point is as good as its posterior allows
at confidence level $\beta$, then pick the best pretend value. Larger
$\beta$ explores harder. The rule comes with regret guarantees under
GP-bandit assumptions, with $\beta$ growing slowly over time so that no
region is starved forever \citep{srinivas2010gaussian}, and it extends
naturally to the time-varying setting that population-based methods
inhabit; PB2, a population-based tuner taught in
Chapter~\ref{sec:population}, uses exactly such a rule, so
the formula is worth remembering.

Noise-aware variants of EI stop pretending the incumbent value
$f^{\star}$ is known, and the reason they must is a bias economists
teach under its own name. With noisy observations the ``best value
seen'' is the \emph{minimum of noisy draws}, and the minimum of noisy
draws is optimistically biased: the trial that reported the best return
in the running example is disproportionately likely to be a decent
configuration that also drew a lucky seed. This is the winner's curse,
and plugging the lucky number into \eqref{eq:ei} as $f^{\star}$
transmits the curse to every later decision: the reported incumbent is
better than the true incumbent, so genuine improvements look smaller
than they are, the acquisition under-rewards them, and the search
stalls defending a score that no configuration, including the
incumbent itself, can reproduce.

The repair is to stop conditioning on
the noisy draws and integrate over what the true values might be. Let
$f(\config_{1:n})$ be the (unknown) true objective values at the
evaluated points; the noisy-EI acquisition is
\begin{equation}
\alpha_{\mathrm{NEI}}(\config)
= \E\!\left[ \max\!\Big( \min_{i \le n} f(\config_i) - f(\config),\; 0
\Big) \right],
\label{eq:nei}
\end{equation}
where the expectation runs over the \emph{joint} posterior of
$f(\config)$ and $f(\config_{1:n})$: the incumbent value inside the
expectation is the true best, not the lucky observed best, and the
average over the posterior prices the luck out.

The integral has no
closed form and is estimated by Monte Carlo, jointly sampling the true
values at old and candidate points; this is the formalization of the
approach used in production at Meta and now standard in BoTorch
\citep{letham2019noisyei, balandat2020botorch}, and its log-stabilized
version inherits the LogEI repair \citep{ament2023logei}.
Reinforcement learning returns and effective-sample-size measurements
are noisy enough that a tuner without the noise-aware acquisition can
waste budget refining estimates of points the GP already knows about
rather than exploring genuinely new candidates, so we ship the
noise-aware variant by default. A user whose observations are nearly
deterministic may switch it off, but RL and sampling workloads keep
it on.

Batch and asynchronous proposals fall out of the same Monte Carlo
formulation, and they answer the wall-clock concern of
Section~\ref{sec:costmodel}. To propose $q$ points jointly, one optimizes
the expected improvement of the \emph{best} of the $q$, which
automatically spreads the batch out (two candidates next to each other
mostly duplicate each other's improvement, so the joint objective prefers
diversity without any explicit diversity penalty).

To propose while $k$
trials are still running, one conditions on their unknown outcomes by
sampling them from the posterior (so-called fantasies) and averages the
resulting decisions \citep{balandat2020botorch}. This is the machinery
that lets model-based search keep a pool of workers busy rather than
idling between sequential suggestions, and it is a hard requirement for
us, not a luxury.

One mechanical detail completes the picture, because ``maximize the
acquisition'' is itself an optimization problem and its failure mode is
the one LogEI fixed. In continuous spaces the standard recipe scores a
few thousand quasi-random candidates (Sobol points again), keeps the
best handful, and polishes each by gradient ascent, which is possible
because $\mu_n$ and $s_n$ in \eqref{eq:gp} are differentiable in
$\config$ and the Monte Carlo acquisitions are built to be
differentiated \citep{balandat2020botorch}. In mixed spaces the
gradient steps apply to the continuous coordinates and local moves to
the discrete ones, the interleaving that Casmopolitan formalizes
(Section~\ref{sec:highdim}). The polish step is exactly where a flat,
underflowed acquisition strands the optimizer at its starting points,
which is why the numerics of the previous subsection are an
engineering requirement rather than a refinement.

\section{Density-ratio and tree-based alternatives}
\label{sec:tpe}

The tree-structured Parzen estimator (TPE) inverts the modeling question,
and the inversion is worth deriving because it explains both TPE's
strengths and its blind spot. Instead of modeling $p(y \mid \config)$ as
the GP does, TPE splits the observed trials at the $\gamma$-quantile of
their scores $y^{\star}$ (say the best 25 percent versus the rest) and
fits kernel density estimates to the two groups of \emph{configurations}:
$\ell(\config)$ for configurations that scored well, $g(\config)$ for the
rest \citep{bergstra2011tpe}. New candidates are proposed where the ratio
$\ell(\config)/g(\config)$ is largest: where good trials cluster and bad
trials are rare (Figure~\ref{fig:tpe-fig}).

The connection to EI is not a
heuristic. Under TPE's model, $p(\config \mid y)$ equals $\ell(\config)$
when $y < y^{\star}$ and $g(\config)$ otherwise, so by Bayes' rule the
marginal is
$p(\config) = \gamma\, \ell(\config) + (1-\gamma)\, g(\config)$, and the
expected improvement below the threshold becomes
\begin{equation}
\alpha_{\mathrm{EI}}(\config)
= \int_{-\infty}^{y^{\star}} (y^{\star} - y)\,
  \frac{p(\config \mid y)\, p(y)}{p(\config)} \, dy
\;=\; \frac{\ell(\config)\int_{-\infty}^{y^{\star}}(y^\star - y)p(y)\,dy}
  {\gamma\, \ell(\config) + (1-\gamma)\, g(\config)}
\;\propto\;
\left( \gamma + (1 - \gamma)\, \frac{g(\config)}{\ell(\config)}
\right)^{-1},
\label{eq:tpe}
\end{equation}
since the integral in the numerator does not depend on $\config$. The
right-hand side is a monotone increasing function of
$\ell(\config)/g(\config)$: maximizing the density ratio \emph{is}
maximizing EI under this model \citep{bergstra2011tpe}.

What the model
gives up is visible in the same derivation: the densities are typically
factorized per coordinate, so the ratio cannot represent ``this learning
rate is good only when the batch size is large,'' an interaction a GP
kernel captures without effort.

\begin{figure}[t]
\centering
\begin{tikzpicture}[scale=1.0, font=\small]
\draw[-{Latex}] (0,0) -- (9,0) node[below left=1mm and 0mm] {learning rate (log scale)};
% good density l(x): concentrated bump
\draw[thick, blue!60!black]
  (0.3,0.05) .. controls (2.2,0.1) and (2.9,0.3) .. (3.5,1.5)
  .. controls (3.8,2.1) and (4.2,2.1) .. (4.5,1.5)
  .. controls (5.1,0.3) and (5.8,0.1) .. (8.6,0.05);
\node[blue!60!black, anchor=south] at (4.0,2.15) {$\ell$: where good trials cluster};
% bad density g(x): broad
\draw[thick, red!60!black, densely dashed]
  (0.3,0.75) .. controls (2.0,0.95) and (3.0,0.55) .. (4.0,0.45)
  .. controls (5.0,0.55) and (6.5,1.0) .. (8.6,0.8);
\node[red!60!black, anchor=west] at (6.1,1.25) {$g$: where the rest landed};
% ratio argmax marker
\draw[-{Latex}, thick] (4.0,-0.55) -- (4.0,-0.08);
\node[anchor=north] at (4.0,-0.6) {propose where $\ell/g$ peaks};
% data ticks
\foreach \x in {3.3,3.7,4.05,4.4,4.7} \draw[blue!60!black, thick] (\x,0.02) -- (\x,0.18);
\foreach \x in {0.8,1.5,2.3,5.6,6.4,7.2,7.9} \draw[red!60!black] (\x,0.02) -- (\x,0.14);
\end{tikzpicture}
\caption{TPE's view of the running example's learning-rate axis
(schematic). Trials are split at a score quantile; the configurations of
the good fraction (solid ticks) get density $\ell$, the rest (dashed
ticks) get $g$, and the next candidate maximizes $\ell/g$, which
by~\eqref{eq:tpe} is the same as maximizing expected improvement under
TPE's model. The construction never represents interactions between axes;
that is the price of its speed and its ease with categorical and
conditional parameters.}
\label{fig:tpe-fig}
\end{figure}

Three properties make TPE the workhorse of Optuna, the open-source
tuning library where sampler development is most active (its record is
audited in Chapter~\ref{sec:libraries}). It
handles conditional (tree-structured) spaces natively, because densities
are fit per branch of the configuration tree; it tolerates categorical
and integer parameters without kernel gymnastics, since a kernel density
over a finite set is just a weighted histogram; and its cost grows gently
with both dimension and trial count, with none of the $O(n^3)$ algebra
of~\eqref{eq:gp}. Its weakness is the flip side of the factorized
densities above: blindness to interactions, and in low dimensions with
smooth objectives a GP squeezes more out of each trial. The recent
tutorial by \citet{watanabe2023tpe} documents the many small design
choices (bandwidths, splitting quantile, prior weight) that separate a
good TPE from a poor one; since Optuna implements these carefully,
wrapping Optuna's sampler is more sensible than reimplementing TPE
ourselves.

SMAC replaces the GP with a random forest surrogate, which natively
digests mixed and conditional spaces and yields uncertainty from the
spread across trees \citep{lindauer2022smac3}. Forest surrogates are less
sample-efficient than GPs on smooth low-dimensional problems but degrade
gracefully on the awkward spaces where GPs struggle, and SMAC3 pairs the
surrogate with multi-fidelity racing (Chapter~\ref{sec:multifidelity}).
We cite it here both as an algorithmic reference point and because it is
the tuner that \citet{wan2022bgpbt} used as their sequential BO baseline.

\section{High-dimensional and mixed spaces}
\label{sec:highdim}

Vanilla GP BO has a reputation for failing beyond ten or twenty
dimensions, which would be disqualifying for the running example's
fifteen. The reputation is only half deserved, and the recent literature
sorts out which half; the sorting matters to us because it changes what
we build first.

The mechanism behind the failures is worth stating before the fixes. In
high dimension, with standard lengthscale priors, the fitted GP tends
toward one of two degenerate postures: lengthscales so short that every
observation is an island (posterior variance is high everywhere, and the
acquisition chases pure exploration into the corners), or a fit that
contorts itself around noise. Each published fix restricts the model in a
different way, and each restriction is a bet about the objective.

TuRBO bets that good local models beat bad global ones. It abandons the
global search problem and maintains local trust regions: boxes around
good points, inside each of which an independent BO runs with its own GP
(Figure~\ref{fig:turbo}). A region's side length grows after runs of
successes (the local model is trustworthy, let it see further) and
shrinks after runs of failures (the local model is misleading, tighten
the box); a region that shrinks below a floor is abandoned and restarted
elsewhere, and a bandit allocates trials across regions
\citep{eriksson2019turbo}. This concentrates samples where the model has
earned trust and scales to hundreds of dimensions; it is the backbone
that BG-PBT adapts for population training (Chapter~\ref{sec:population}),
which is the specific reason this document cares about its details.

\begin{figure}[t]
\centering
\begin{tikzpicture}[scale=0.9, font=\small]
\draw[thick] (0,0) rectangle (5.2,4.2);
\node[anchor=south west, font=\small\itshape] at (0.05,4.22) {search space (2D slice)};
% incumbent
\fill[blue!70] (2.4,2.0) circle (2.5pt);
% trust region boxes
\draw[blue!60!black, thick] (1.7,1.3) rectangle (3.1,2.7);
\draw[blue!60!black, densely dashed] (1.0,0.6) rectangle (3.8,3.4);
\draw[blue!60!black, dotted, thick] (2.05,1.65) rectangle (2.75,2.35);
% arrows and labels
\draw[-{Latex}] (5.6,3.4) -- (3.85,3.1);
\node[anchor=west, align=left] at (5.6,3.4) {grow after successes:\\the local model has\\earned a longer leash};
\draw[-{Latex}] (5.6,1.0) -- (2.78,1.85);
\node[anchor=west, align=left] at (5.6,1.0) {shrink after failures;\\below a floor, restart\\elsewhere with a fresh model};
\draw[-{Latex}, blue!70] (2.4,3.62) -- (2.4,2.12);
\node[anchor=south, blue!70] at (2.4,3.66) {incumbent};
\end{tikzpicture}
\caption{TuRBO's trust region around the incumbent (schematic). All
modeling and acquisition happen inside the current box; success and
failure counters resize it, and collapse triggers a restart. The bet is
that a GP asked only local questions stays reliable in dimensions where
a global GP degenerates, and the same box-resizing logic reappears
inside BG-PBT in Chapter~\ref{sec:population}.}
\label{fig:turbo}
\end{figure}

SAASBO (sparse axis-aligned subspace BO) bets that only a few axes
matter. It keeps the global view but
places sparsity-inducing half-Cauchy priors on the kernel's inverse
squared lengthscales, so that every axis starts effectively switched off
(infinite lengthscale) and only axes with evidence of mattering wake up;
inference is fully Bayesian over these priors, which buys striking sample
efficiency in hundreds of nominal dimensions and costs enough per
iteration to confine the method to small evaluation budgets
\citep{eriksson2021saasbo}. If the effective-dimension picture behind
Figure~\ref{fig:gridrandom} is right for a problem, SAASBO is that
picture taken seriously as a prior. A related line grows a random
embedding of the space as evidence accumulates rather than committing to
one up front, with a mixed-space successor for combinatorial problems
\citep{papenmeier2022baxus, papenmeier2023bounce}.

Then the twist that reorders the whole subsection. \citet{hvarfner2024vanilla}
showed that much of the classical failure was a prior mistake rather than
a fundamental one: simply scaling the GP's lengthscale prior with the
dimension, so that the model's complexity stays controlled as axes are
added, lets plain BO with standard acquisitions outperform the
specialized high-dimensional methods on common benchmarks, with no
structural assumption about the objective. The intuition connects
straight back to \eqref{eq:mll}: the complexity penalty a kernel pays
scales with how much the function can wiggle across the whole space,
which grows with dimension unless lengthscales grow too, so the fix is a
prior that expects longer lengthscales in higher dimension. The
practical ordering for us follows: fix the priors first (cheap, one line
of configuration), reach for trust regions when the space is genuinely
large or the budget is generous, and reserve SAASBO for sample-starved,
high-stakes searches.

Mixed spaces need kernels that understand their coordinates. Continuous
axes take standard kernels; categorical axes take overlap or
Hamming-style kernels (two configurations are similar to the extent they
agree, with no notion of ``between''); ordinal axes (integers with
order, like batch size or network width) should be treated as coarse
continuous axes rather than unordered categories, because width 256
really is between width 128 and width 512, and throwing that ordering
away discards exactly the structure a kernel can use.

Casmopolitan
combines these with trust regions and interleaved acquisition
optimization, alternating local search over the discrete coordinates
with gradient steps over the continuous ones \citep{wan2021casmopolitan},
and the BG-PBT paper reports that explicitly respecting ordinal
structure gives measurable gains in the RL setting \citep{wan2022bgpbt}.

An alternative with BoTorch support is probabilistic reparameterization,
which relaxes discrete choices into distributions over them and
optimizes the relaxation with gradients \citep{daulton2022pr}. Our
product needs exactly one good mixed-space GP searcher; Casmopolitan-style
kernels inside a trust region, reusing BoTorch primitives, is the path of
least resistance and doubles as the core of a future BG-PBT
implementation.

\section{The evidence, and how it was produced}
\label{sec:bo-evidence}

Everything so far argues that model-based search \emph{should} work.
A skeptical reader is entitled to ask what would count as evidence that
it does, and to ask it before accepting any verdict, so this section
describes the experiments themselves, not just their conclusions.

The evidence problem in this field is selection bias, and it is worth
naming before weighing any result. A paper introducing a method chooses
its own benchmarks, its own baselines, and its own tuning of those
baselines, and only successful comparisons reach print; a reader
cannot tell how much of a reported gain is method and how much is
choice of arena. The strongest evidence therefore comes from settings
the method's authors did not control. Two such settings anchor this
chapter's verdicts.

The first is the NeurIPS 2020 black-box optimization challenge, a
public competition in which entrant tuners were run on real
machine-learning hyperparameter tuning problems at matched evaluation
budgets, with the organizers rather than the entrants choosing the
tasks. The organizers' published analysis found the Bayesian
optimization entries clearly ahead of random search across the board
\citep{turner2021bayesian}: an adversarial arena, fixed budgets, and
the baseline of Chapter~\ref{sec:modelfree} still cleanly beaten. The
winning entry, HEBO, is instructive for \emph{why} it won. Its
authors' paper opens not with a new acquisition function but with
hypothesis tests on 108 real tuning tasks showing that
heteroscedastic, non-stationary objectives are the rule rather than
the exception, and the winning design is essentially the standard GP
machinery of this chapter plus the input and output transformations
that make such objectives look well behaved to it
\citep{cowenrivers2022hebo}. The reading we take: the competition
validates the framework, and the margin came from respecting the
data's geometry, which is an engineering discipline we can copy, not a
secret.

The second setting is the controlled RL tuning study of
\citet{eimer2023hyperparameters}. Three algorithms (PPO, SAC, DQN)
across eleven environments; tuners compared at matched budgets of ten
to sixty-four full training runs, which is the regime our routine
studies actually inhabit; candidates evaluated on several tuning seeds
with final results reported on disjoint test seeds the tuner never
saw; and the tuners' own overhead measured. The finding that matters
here: sweeping hyperparameters across eleven environments, the fitted
response surfaces looked smooth and well behaved, with the difficulty
of RL tuning coming overwhelmingly from seed noise rather than from a
pathological objective. Smooth surfaces are exactly what the kernel
assumption \eqref{eq:matern} posits, and severe observation noise is
exactly what the noisy-EI machinery \eqref{eq:nei} prices in, so the
one controlled study of our hardest workload lands on both of this
chapter's design choices. (The study's full design and its
budget-regime findings occupy Chapter~\ref{sec:population}, where the
population line's evidence lives.)

Two honest boundaries on this evidence. It does not rank GP against
TPE universally: which surrogate wins depends on the space's shape and
the budget, which is why the package carries both and why the
comparison that matters will be run on our own pinned tasks. And the
LogEI results, showing the repaired classic acquisitions matching or
beating the more elaborate ones across standard benchmarks
\citep{ament2023logei}, are the authors' own experiments; we treat
them as strong because the claim is negative and numerical (a bug
explained away a literature's worth of differences) and because the
mechanism, the underflow arithmetic of
Section~\ref{sec:acquisition}, can be checked from first principles
on one page, which we did.

\section{The recommendation, argued from four angles}
\label{sec:bo-angles}

The chapter's verdict is that a carefully configured GP searcher with
log-stabilized, noise-aware acquisitions is the core of the package,
with TPE wrapped alongside for conditional spaces. Four independent
lines of argument arrive at that verdict, and a reader should demand
that they agree, because each could fail separately.

From decision theory: when a trial costs GPU-hours, the value of
deliberation between trials is high, and \eqref{eq:ei} is literally
the option value of the next experiment under current beliefs. Any
method that ignores posterior uncertainty is discarding the one
quantity that says where deliberation pays.

From statistics: the machinery is not exotic. The surrogate is
kriging, the noise handling is ridge shrinkage with an estimated
regularizer, the winner's-curse correction is the same bias argument
auction theory makes. A reader who trusts those tools separately is
most of the way to trusting their composition; nothing in the chapter
asks for a new inferential principle.

From evidence: an author-independent competition on real tuning tasks
put the framework clearly ahead of the honest baseline, the winner's
margin traced to checkable preprocessing rather than proprietary
insight, and the one controlled study of our own hardest workload
found exactly the surface smoothness and noise structure this design
assumes.

From operations: every component the chapter recommends exists,
maintained, in BoTorch and Optuna (Chapter~\ref{sec:libraries}), so
the build cost is composition rather than research; batch and
asynchronous proposals keep a cluster busy; and the known numerical
failure mode has a published, adopted repair.

And the check on all four: the verdict is falsifiable in our own
validation suite. The GP searcher must beat Sobol sampling at matched
budgets on the pinned tasks of Chapter~\ref{sec:roadmap}, with the
comparison repeated across seeds and reported with uncertainty. If it
cannot, the verdict, not the suite, gives way; that is the standard we
would demand of anyone else's recommendation, so the document applies
it to its own.

\begin{decisions}
\noindent\textbf{Decisions from this chapter.} The choice among
surrogates comes down to four manager-checkable questions: how few
trials before the model helps, what kinds of knobs it digests, how it
copes with noise, and what it costs us to ship.
\begin{itemize}[leftmargin=1.3em, itemsep=2pt]
\item \textbf{GP searcher with qLogEI on BoTorch} (include; the core
build, weeks). Extracts the most from the fewest trials on mostly
continuous spaces of up to a few dozen dimensions, handles noise
honestly via the noisy-EI machinery, proposes batches so workers never
idle. This is the searcher for expensive trials, which is to say for
the studies where the tuner's quality is actually worth money. The
challenge evidence is on its side: model-based entries beat random
search across the board, and the winner was a GP method
\citep{turner2021bayesian, cowenrivers2022hebo}.
\item \textbf{TPE by wrapping Optuna} (include; days). Useful from a
handful of trials, natively handles conditional and categorical spaces
that GPs find awkward, and costs us almost nothing because we wrap the
best-maintained implementation rather than writing one. Its known
blind spot, interactions between knobs, is the price
(equation~\ref{eq:tpe}); for smooth low-dimensional problems the GP
squeezes more from each trial.
\item \textbf{Raw EI, anywhere} (exclude). The underflow analysis is
not a nicety: past a modest incumbent quality the naive formula returns
exactly zero over most of the space and the optimizer stalls
\citep{ament2023logei}. Log variants only.
\item \textbf{Entropy-search acquisitions} (exclude). Elegant, costly
to implement well, and the controlled comparisons show the repaired EI
family matching or beating them \citep{ament2023logei}. No capability
we need survives their exclusion.
\item \textbf{Dimension-scaled lengthscale priors} (include;
one configuration line). The cheapest high-dimensional fix on record
\citep{hvarfner2024vanilla}. Always on.
\item \textbf{Trust-region searcher with mixed-space kernels}
(include; weeks). Justified twice over: it is the credible
high-dimensional searcher, and it is a prerequisite we would build
anyway for BG-PBT in Chapter~\ref{sec:population}.
\item \textbf{SAASBO} (defer). Striking sample efficiency when trials
are precious and axes are few-but-unknown, at a per-iteration compute
cost that only suits small, high-stakes studies. Revisit when such a
study appears.
\item \textbf{SMAC} (do not adopt as a dependency). We use it as a
correctness oracle in validation; its random-forest surrogate niche
(huge conditional spaces) is not where our workloads live.
\end{itemize}
\end{decisions}

\section{What would overturn these verdicts}

The GP searcher's tier-0 status is tied to V04: if it cannot beat
the Sobol floor on our pinned tasks, the complexity is not justified
and it is demoted to an option. TPE's wrapping strategy holds as
long as Optuna's maintenance remains current (audited in
Chapter~\ref{sec:libraries}); if Optuna's TPE implementation
stagnates or breaks backward compatibility repeatedly, the verdict
changes to a native build. The trust-region searcher's shared role
(high-dimensional search and population explore) is load-bearing; if
the population line is demoted by V08, the trust-region design is
re-argued on single-objective merits alone. Raw EI and
entropy-search exclusions are reopened if a published study on RL or
sampler workloads shows either beating qLogEI at matched budgets, at
which point the evidence section is updated and the box reconsidered.
